\setcounter{section}{0}

\chapter*{Trouver un sujet}
\subsection*{\centering Objectif : Choisir son sujet.}

Un bon sujet pour le Grand Oral est avant tout un sujet qui passionne l’élève. Plus l’élève est intéressé par ce qu’il présente, et plus il lui sera facile de préparer son oral et d’être à l’aise avec son discours.\\

\section{Exemples de sujets (Ministère de l'éducation)}
\textbf{L’histoire de l’informatique}
\begin{itemize}
    \item Femmes et numérique : quelle histoire ? quel avenir ?
    \item Ada Lovelace, pionnière du langage informatique
    \item Alan Turing, et l’informatique fut
    \item Quelle est la différence entre le web 1.0 et le web 2.0 ?
\end{itemize}

\textbf{Langages et programmation}
\begin{itemize}
\item P = NP, un problème à un million de dollars ?
\item Tours de Hanoï : plus qu’un jeu d’enfants ?
\item Les fractales : informatique et mathématiques imitent-elles la nature ?
\item De la récurrence à la récursivité
\item Les bugs : bête noire des développeurs ?
\item Comment rendre l’informatique plus sûre ?
\end{itemize}

\textbf{Données structurées et structures de données}
\begin{itemize}
\item L’informatisation des métros : progrès ou outil de surveillance ?
\item Musique et informatique : une alliance possible de l’art et de la science ?
\end{itemize}

\textbf{Algorithmique}
\begin{itemize}
    \item Comment créer une machine intelligente ?
    \item Comment lutter contre les biais algorithmiques ?
    \item Quels sont les enjeux de la reconnaissance faciale (notamment éthiques) ?
    \item Quels sont les enjeux de l’intelligence artificielle ?
    \item Transformation d’images : Deep Fakes, une arme de désinformation massive ? La fin
    de la preuve par l’image ?
    \item Qu’apporte la récursivité dans un algorithme ?
    \item Quel est l’impact de la complexité d’un algorithme sur son efficacité ?
\end{itemize}

\textbf{Bases de données}
\begin{itemize}
    \item Données personnelles : la vie privée en voie d’extinction ?
    \item Comment optimiser les données ?
\end{itemize}

\textbf{Architectures matérielles, systèmes d’exploitation et réseaux}
\begin{itemize}
    \item L’ordinateur quantique : nouvelle révolution informatique ?
    \item La course à l’infiniment petit : jusqu’où ?
    \item Peut-on vraiment sécuriser les communications ?
    \item Quelle est l’utilité des protocoles pour l’internet ?
    \item Cyberguerre : la 3ème guerre mondiale ?
\end{itemize}

\textbf{Interfaces Hommes-Machines (IHM)}
\begin{itemize}
    \item Smart cities, smart control ?
    \item La réalité virtuelle : un nouveau monde ?
    \item La voiture autonome, quels enjeux ?
\end{itemize}

\textbf{Impact sociétal et éthique de l’informatique}
\begin{itemize}
    \item Comment protéger les données numériques sur les réseaux sociaux ?
    \item Quelle est l’empreinte carbone du numérique en termes de consommation ?
    \item Pourquoi chiffrer ses communications ?
    \item Les réseaux sociaux sont-ils compatibles avec la politique ?
    \item Les réseaux sociaux sont-ils compatibles avec le journalisme ?
    \item Les réseaux sociaux permettent-ils de lutter contre les infox ?
    \item L’informatique va-t-elle révolutionner le dessin animé ?
    \item L’informatique va-t-elle révolutionner la composition musicale ?
    \item L’informatique va-t-elle révolutionner l’art ?
    \item L’informatique va-t-elle révolutionner le cinéma ?
    \item L’informatique va-t-elle révolutionner la médecine ?
    \item L’informatique va-t-elle révolutionner la physique ?
    \item L’informatique va-t-elle révolutionner l’entreprise ?
    \item Le numérique : facteur de démocratisation ou de fractures sociales ?
    \item Informatique : quel impact sur le climat ?
\end{itemize}



\section{D'autres sujets par des Youtubeurs}
Proposition de vidéos de vulgarisation portant sur différentes thématiques.
\begin{itemize}
    \item Pixees Scienceparticipative : Un condensé d'histoire de l'informatique
    \item Pixees Scienceparticipative : La notion d'Information (Gérard Berry)
    \item Heu?reka : Traders à la vitesse de la lumière - Heu?reka \#9
    \item Heu?reka : Trader à la vitesse de la lumière (le retour!) - Heu?reka \#12
    \item Heu?reka : Crypto-investissements : éviter les pièges - Heu?reka \& AMF
    \item Le dessous des cartes : câbles fibres sous-marins
    \item Sciences4all : la démocratie sous l’angle de la théorie des jeux 
    \item Le Réveilleurs : les terres rares 
    \item Fouloscopie : peut-on prédire les futurs crimes 
    \item Avides de recherches : l’amour en ligne 
    \item SciencesEtonnante : la fourmie de langton
    \item ScienceEtonnante : Comment ces IA inventent-elles des images ?
    \item ScienceEtonnante : Comment les I.A. font-elles pour comprendre notre langue ?
    \item ScienceEtonnante : Les codes secrets
    \item ScienceEtonnante : JE CRAQUE WORDLE !(grâce aux maths)
    \item Pass sciences : jeu de la vie
    \item Scilabus : Ada lovelace
    \item Mickael launay : un bug géométrique dans google map
    \item Hacker dashery : P vs NP
    \item Math adultes : blockchain comment ça marche
    \item Vortex : fake news tous coupables 
    \item Vortex : réseaux sociaux une nouvelle facon de devenir immortel 
    \item Zeste De Science : au corps à corps avec les robots
    \item StupidsEconomics : Notre argent, la youtube monnaie 
    \item Pause Process : 48 le hack
    \item Chat Sceptique : l'erreur excel qui a couté la vie à 1500 personnes
    \item Chat Sceptique : les algorithmes révolutionnaires racontés à mon chat
    \item e-penser : Alan Turing - Enigma, ordinateur et pomme empoisonnée
    \item Etincelles, EPFL : Réseaux de neuronnes
    \item Balade mentale : Trouver la DATE DE VOTRE MORT ? Bibliothèque de BABEL / Paradoxe du singe savant
    \item Balade mentale : Ada lovelace \& l'ORDINATEUR MÉCANIQUE
    \item Monsieur Phi : ChatGPT : la philosophie du baratin
    \item Monsieur Phi : La preuve définitive que ChatGPT ne comprend rien
    \item Monsieur Phi : ChatGPT peut-il passer le bac de philo ?
    \item Dans le secret du fonctionnement des IA - Science En Questions : Dans le secret du fonctionnement des IA - Science En Questions. L’Explicabilité de l’IA : pourquoi est-ce crucial ?
\end{itemize}

\section{Encore d'autres sujets possibles}
\begin{itemize}
    \item Pourquoi Javascript est-il devenu le langage le plus utilisé au monde ?
    \item Pourquoi chiffrer les communications ?
    \item L'informatique va-t-elle révolutionner le dessin animé ?
    \item Comment l'informatique a-t-elle révolutionné la vie des gens dans le domaine du divertissement au fil des années ?
    \item Comment fonctionnement les bases de données distribuées ?
    \item Nos informations personnelles sont-elles en sécurité dans les bases de données ?
    \item Comment l'informatique s'inspire du vivant ?
    \item Enigma et décryptage. En quoi l'informatique a-t-elle fait évoluer les maths ?
    \item Comment évolue la cybersécurité ?
    \item Modélisation 3D et graphes.
    \item L'art algorithmique est-il vraiment un art ?
    \item Comment protéger les données numériques sur les réseaux
    \item Quels sont les enjeux des cyberguerres ?
    \item Fractales, l'info et les maths imitent-elles la nature ?
    \item Les réseaux sociaux sont-ils compatibles avec la politique ?
    \item Avenir de la réalité virtuelle ?
    \item Les réseaux sociaux sont-ils compatibles avec le journalisme ?
    \item En quoi la cryptologie a-t-elle permis de raccourcir la Deuxième Guerre mondiale ?
    \item Comment fonctionne Tor ?
    \item Comment flash est-il passé du statut d'utilisable à mort en 2021 ?
    \item Pourquoi Python a-t-il été inventé ?
    \item Comment l'automate cellulaire peut-il imiter les êtres vivants par de simples règles informatiques ?
    \item Comment HTML 5 a révolutionné le développement web ?
    \item Quelles ont été les évolutions de l'utilisation d'internet ?
\end{itemize}

\newpage

\textbf{Faites des recherches puis entourer 5 sujets susceptibles de vous intéresser}, pour chacun de ces sujets, vous devez en quelques phrases expliquer pourquoi il vous intéresse de manière convaincante.\\

\noindent Sujet 1 :
\vspace{1cm}\\
Sujet 2 :
\vspace{1cm}\\
Sujet 3 :
\vspace{1cm}\\
Sujet 4 :
\vspace{1cm}\\
Sujet 5 :
\vspace{1cm}

\section{Mes intérêts}

Énumérer cinq de vos centres d'intérêt actuels (loisirs, actualité, sciences, société, etc.), puis notez trois mots-clés pour votre deuxième spécialité (qu'est-ce qui vous intéresse dans cette spécialité).
L'idée est de créer des ponts.\\
Exemple :
\begin{itemize}
    \item Spé 2 (Maths) $\rightarrow$ Mots-clés : Modélisation, IA, Cryptographie.
    \item Intérêt personnel $\rightarrow$ Jeux vidéos.
    \item Pont : Comment l'IA est-elle utilisée dans la création de personnages non-joueurs (PNJ)?
\end{itemize}
\vspace{0.5cm}
\noindent Centres d'intérêts :
\vspace{1cm}\\
Mots-clés :
\vspace{1cm}\\
Pont :
\vspace{1cm}

\section{De l'intérêt à la question}

 Sur la base des thèmes entourés et des ponts trouvés, écrire trois questions de Grand Oral potentielles.\\

Rappel : Une bonne question est souvent un problème qui ouvre le débat (Exemple: Quels sont les enjeux de la reconnaissance faciale (notamment éthiques).\\

\noindent Question 1 :
\vspace{1cm}\\
Question 2 :
\vspace{1cm}\\
Question 3 :
\vspace{1cm}